\section{Lean formalization plan}
\label{sec:leanplan}

\subsection{Pinned upstream API}

The outline imports only checked repository modules and uses their public
objects:
\[
\code{PR},\quad \code{PTm},\quad \code{PRAAx},\quad
\code{Prf},\quad \code{PRA.qfInd},\quad
\code{NFm.ofPRA},\quad \code{pra_iff_native}.
\]
The essential proof combinators are
\code{Prf.subst}, \code{Prf.weaken}, \code{Prf.trans},
\code{Prf.impI}, \code{Prf.impE}, \code{Prf.allI},
\code{Prf.allE}, \code{Prf.eqSubst1}, \code{Prf.eqSymm}, and
\code{Prf.eqCongr}.  Primitive-recursion equations come from
\code{PRA.precZero} and \code{PRA.precSucc}.

\subsection{File dependency graph}

\begin{enumerate}
\item \code{ObjectTraceCompiler.lean}: formula builders, trace system, run
equations, four proof milestones, main compiler.
\item \code{BoundedFormulaCompiler.lean}: Boolean symbols, folds, fixed bounded
formula syntax, reflection theorem.
\item \code{RelationalWitnesses.lean}: finite tower codes, viability, least
path, stable observation.
\item \code{CertifiedProgram.lean}: finite program/DAG syntax, reverse demands,
forward compilation to a trace system.
\item \code{ApplicationTargets.lean}: guarded reciprocal, height/list, cubic
bisection, carrier, and affine-gluing targets.
\end{enumerate}

The first file can be completed independently and should be merged before the
others.  It is the narrowest bridge to the checked backend.

\subsection{Exact trace layouts}

The object trace file uses these de Bruijn layouts:
\begin{center}
\begin{tabular}{lll}
\toprule
Formula & Arity & Variables by index \\
\midrule
\code{BaseFm} & 1 & \([u]\) \\
\code{StepFm} & 3 & \([i,s,u]\) \\
\code{FinalFm} & 2 & \([s,u]\) \\
\code{InvRunFm} & 2 & \([i,u]\) \\
\code{ResultFm} & 0 & closed, with one universal binder \\
\bottomrule
\end{tabular}
\end{center}
Term builders \code{app1}, \code{app2}, and \code{app3} make argument order
visible and reduce substitution errors.

\subsection{Milestone 1: base case}

The target is
\begin{lstlisting}[language=Lean]
theorem trace_base (T : PRTraceSystem) (h : TraceProofs T) :
  [GoodFm T] |- [PRAAx] (InvRunFm T).subst (inst zero)
\end{lstlisting}
The proof steps are:
\begin{enumerate}
\item instantiate \code{h.base} at \code{var 0};
\item weaken it to context \code{[GoodFm T]};
\item eliminate the implication using \code{Prf.hyp0};
\item obtain \code{run_zero T (var 0)};
\item transport the state argument with \code{Prf.eqSubst1};
\item close substitution equalities by \code{simp} lemmas for \code{app1/2/3}.
\end{enumerate}
A useful helper is a congruence lemma for \code{IsTrue (app3 T.inv i - u)}.

\subsection{Milestone 2: successor case}

The target is exactly the successor premise expected by \code{PRA.qfInd}:
\begin{lstlisting}[language=Lean]
theorem trace_step (T : PRTraceSystem) (h : TraceProofs T) :
  InvRunFm T :: [GoodFm T].map Fm.shift |- [PRAAx]
    (InvRunFm T).subst (repl (succ (var 0)))
\end{lstlisting}
Instantiate \code{h.step} with
\((i,\Run(i,u),u)\), eliminate implications using \code{Prf.hyp0} and
\code{Prf.hyp1}, and transport along \code{run_succ}.  Most work is normalizing
the shifted guard and the simultaneous substitution.

\subsection{Milestone 3: induction and final contract}

The induction theorem is already one line:
\begin{lstlisting}[language=Lean]
exact PRA.qfInd (InvRunFm T) (invRun_qf T)
  (trace_base T h) (trace_step T h)
\end{lstlisting}
Then use
\begin{lstlisting}[language=Lean]
Prf.allE (app1 T.bound (var 0)) (trace_all T h)
\end{lstlisting}
to obtain the invariant at the terminal index.  Instantiate \code{h.final} with
the terminal run and input, eliminate implications, discharge the guard, and
close the input binder.

\subsection{Milestone 4: bounded Boolean folds}

Introduce canonical Boolean-valued symbols and equations.  For each fixed body
symbol \(f\), define \code{allFold f}, \code{anyFold f}, and
\code{leastFold f} by \code{PR.prec}.  Prove:
\begin{itemize}
\item zero and successor equations from \code{PRA.precZero/precSucc};
\item Boolean range invariants;
\item reflection against bounded formula abbreviations;
\item least-witness soundness and minimality.
\end{itemize}
Use a fixed family indexed at the Lean level; do not define a universal PR
evaluator internally.

\subsection{Milestone 5: relational viability}

Encode a tower by fixed symbols for level bounds, candidate validity, and
edges.  Define viability by a backward recursion on the distance to horizon.
The principal object theorem is a Boolean equivalence between viability and a
bounded path verifier.  A second theorem proves correctness of least viable
successor.  Both use quantifier-free induction on the finite distance.

\subsection{Milestone 6: certified program compilation}

Replace semantic \code{Prop} fields in the program outline by:
\begin{itemize}
\item formula builders over \code{PTm};
\item local derivation fields;
\item fixed arity and state-layout records;
\item a compiler producing \code{PRTraceSystem} and \code{TraceProofs}.
\end{itemize}
Sequential composition uses \code{Prf.trans}.  Bounded loops invoke the trace
compiler recursively.  A static DAG is first lowered to reverse-demand and
forward-table loops, then compiled once.

\subsection{Suggested build discipline}

\begin{enumerate}
\item Place the modules under the root Lake workspace, in namespace
\code{LeanProofs.ConciseIPRA.}\allowbreak\code{FinitaryInterface}.
\item Add a lint command rejecting \code{sorryAx}, explicit \code{axiom}, and
unintended imports.
\item Build after each milestone; keep formula-normalization lemmas separate
from mathematical proofs.
\item Add tiny closed examples: identity trace, addition fold, least witness,
three-level viability tower, and cubic bisection step.
\item Only after the generic modules are checked should large packets import
them.
\end{enumerate}

\subsection{Estimated residual risk}

The trace compiler is low mathematical risk and moderate de Bruijn-engineering
risk.  Boolean folds are moderate risk because a clean fixed-family API must be
chosen.  Relational viability is low mathematical risk but may produce large
formula terms.  The program compiler is architectural rather than logical;
its main risk is over-generalization.  It should initially support only static
DAGs and bounded loops, adding richer syntax after concrete migrations expose
a need.
